\subsection{সম্পর্কিত Table থেকে Data বের করা}

অংশ~\ref{sec:ch06-join}-এ JOIN-এর Syntax ও রেকর্ড মেলানো দেখানো হয়েছে।
এখানে সেই পদ্ধতি Department--Student-এর 1:M এবং পরে Junction Table-এর
M:N সম্পর্কে প্রয়োগ করি। সম্পর্ক সংযোগের নিয়ম ঠিক করে; কুয়েরি সেই
সংযোগ ব্যবহার করে প্রয়োজনীয় ফল বের করে।

\noindent\begin{minipage}{\linewidth}
\chSixExample{১০.১}{Science বিভাগের শিক্ষার্থীদের StudentID ও Name দেখাও।}

\noindent
\begin{minipage}[c]{0.58\linewidth}

\begin{verbatim}
SELECT Student.StudentID,
       Student.Name
FROM Student, Department
WHERE Student.DeptID =
      Department.DeptID
  AND Department.DeptName =
      'Science'
ORDER BY Student.StudentID;
\end{verbatim}

\end{minipage}
\hfill
\begin{minipage}[c]{0.04\linewidth}
\raggedright
$\Rightarrow$
\end{minipage}
\hfill
\begin{minipage}[c]{0.34\linewidth}
\raggedright
\footnotesize

\textbf{ফল}

\smallskip

\begin{tabular}{|l|l|}
\hline
\chSixHeader
StudentID & Name\\
\hline
102 & Rahim\\\hline
109 & Nabila\\\hline
150 & Mim\\
\hline
\end{tabular}

\end{minipage}
\end{minipage}\par\smallskip


\noindent
\textbf{ফলটি কীভাবে পেলাম?} প্রথমে \texttt{Student.DeptID = Department.DeptID}
দিয়ে প্রতিটি শিক্ষার্থীর সঙ্গে তার নিজের Department মেলে। যেমন Rahim-এর
D01 মেলে Science-এর D01-এর সঙ্গে; তাঁর সঙ্গে D02 বা D03 মেলে না।
এরপর \texttt{DeptName='Science'} শর্তে D01-এর Rahim, Nabila ও Mim থাকে।
SELECT শুধু StudentID ও Name দেখায়; ORDER BY StudentID-এর ঊর্ধ্বক্রম নিশ্চিত করে।
শুধু Science-এর শর্ত লিখে Key মেলানো বাদ দিলে অন্য বিভাগের শিক্ষার্থীও
Science সারির সঙ্গে যুক্ত হয়ে ভুল ফল দিতে পারে।


% ============================================================
% MANY TO MANY
% ============================================================

